\subsection{Partitions of the character space}

PROPOSITION 13. --- Let A be a commutative unital Banach algebra. Let U₁ and U₂ be open subsets of X(A) forming a partition of X(A). Then there exists a unique idempotent j in A such that the Gelfand transform G(j) equals 1 on U₁ and 0 on U₂.

Let us identify the space \(\mathsf{X}(\mathsf{A})\) with a compact subset of \(\mathbf{C}^{\mathrm{A}}\) by the map \(\chi \mapsto (\chi (a))_{a\in \mathrm{A}}\) (cf. no. 6 of I, p. 6 and cor. of th. 1 of I, p. 29). The subsets \(\mathrm{U}_1\) and \(\mathrm{U}_2\) of the uniform space \(\mathbf{C}^{\mathrm{A}}\) are compact and disjoint.

By TG, II, p.~31, prop. 4, there exists a finite subset M of A and disjoint open sets \(\mathrm{V}_{1}\) and \(\mathrm{V}_{2}\) of \(\mathbf{C}^{\mathrm{M}}\) such that

\[
p (\mathrm{U} _ {1}) \subset \mathrm{V} _ {1}, \qquad p (\mathrm{U} _ {2}) \subset \mathrm{V} _ {2},
\]

where \(p\) is the canonical projection of \(\mathbf{C}^{\mathrm{A}}\) on \(\mathbf{C}^{\mathrm{M}}\) .

Let \(a_1, \ldots, a_n\) be the distinct elements of M, and identify \(\mathbf{C}^{\mathrm{M}}\) with \(\mathbf{C}^n\) . We have \(\mathrm{Sp}_{\mathrm{A}}^{n}(a_{1}, \ldots, a_{n}) \subset p(\mathsf{X}(\mathrm{A})) \subset \mathrm{V}_{1} \cup \mathrm{V}_{2}\) since \(\mathrm{U}_{1} \cup \mathrm{U}_{2} = \mathsf{X}(\mathrm{A})\) . Let \(f\) be the function on \(\mathrm{V}_{1} \cup \mathrm{V}_{2}\) equal to 1 on \(\mathrm{V}_{1}\) and to 0 on \(\mathrm{V}_{2}\) . We have \(f \in \mathcal{O}(\mathrm{V}_{1} \cup \mathrm{V}_{2})\) . Set \(j = f(a_{1}, \ldots, a_{n})\) . Since \(f^{2} = f\) , one has \(j^{2} = j\) . By cor. 1 of I, p.~66, we have \(\chi(j) = 1\) if \(\chi \in \mathrm{U}_{1}\) and \(\chi(j) = 0\) if \(\chi \in \mathrm{U}_{2}\) , which proves the existence of the required idempotent.

On the other hand, if \(j_{1}\) is an idempotent of A, the relations \(j^{2} = j\) and \(j_{1}^{2} = j_{1}\) imply \((j - j_{1})(j + j_{1} - 1) = 0\) . If \(\mathcal{G}(j_{1}) = \mathcal{G}(j)\) , the Gelfand transform of \(j + j_{1} - 1\) takes values in \(\{-1, 1\}\) , hence \(j + j_{1} - 1\) is invertible (prop. 6 of I, p.~37), whence \(j = j_{1}\) .

COROLLARY. --- Let A be a commutative unital Banach algebra. The following assertions are equivalent:

\begin{enumerate}
\def\labelenumi{\alph{enumi})}
\item
  The space of characters X(A) is not connected;
\item
  There exists an idempotent element of A different from 0 and 1.
\item
  The algebra A is isomorphic to the product of two nonzero Banach algebras.
\end{enumerate}

The proposition demonstrates that \(a\) ) implies \(b\) ). If \(j\) is an idempotent of A, let \(\mathrm{I}_1 = j\mathrm{A}\) and \(\mathrm{I}_2 = (1 - j)\mathrm{A}\) . Then \(\mathrm{I}_1\) and \(\mathrm{I}_2\) are closed ideals of A, and \(\mathrm{I}_1 + \mathrm{I}_2 = \mathrm{A}\) . If \(j \notin \{0,1\}\) , the ideals \(\mathrm{I}_1\) and \(\mathrm{I}_2\) are distinct from A. On the other hand, the ideal \(\mathrm{I}_1\) (resp. \(\mathrm{I}_2\) ) is the set of elements \(x\) of A such that \(jx = x\) (resp. \((1 - j)x = x\) ), hence \(\mathrm{I}_1 \cap \mathrm{I}_2 = \{0\}\) . The algebra A is then identified with the product \(\mathrm{A} / \mathrm{I}_1 \times \mathrm{A} / \mathrm{I}_2\) . Thus, the assertion \(b\) ) implies \(c\) ). Finally, if A is isomorphic to \(\mathrm{A}_1 \times \mathrm{A}_2\) , the space \(\mathsf{X}(\mathsf{A})\) is identified with the sum space of \(\mathsf{X}(\mathsf{A}_1)\) and of \(\mathsf{X}(\mathsf{A}_2)\) (I, p.~6, no. 6), hence \(c\) ) implies \(a\) ).

PROPOSITION 14. --- Let A be a commutative Banach algebra without radical. For A to admit an identity element, it is necessary and sufficient that X(A) be compact.

The condition is necessary (I, p.~29, corollary). Suppose \(\mathsf{X}(\mathsf{A})\) compact. Let \(\widetilde{\mathsf{A}}\) be the Banach algebra obtained from A by adjoining an identity element, and identify \(\mathsf{X}'(\mathsf{A})\) with \(\mathsf{X}(\widetilde{\mathsf{A}})\) . The complement of \(\mathsf{X}(\mathsf{A})\) in \(\mathsf{X}(\widetilde{\mathsf{A}})\) is reduced to the character \(\chi_0\) of \(\widetilde{\mathsf{A}}\) whose kernel is A. The sets \(\mathsf{X}(\mathsf{A})\) and \(\{\chi_0\}\) are open in \(\mathsf{X}(\widetilde{\mathsf{A}})\) . By prop. 13, there exists an element \(j \in \mathsf{A}\) such that \(\chi(j) = 1\) for \(\chi \in \mathsf{X}(\mathsf{A})\) , and \(\chi_0(j) = 0\) . We therefore have \(j \in \mathsf{A}\) .

Let then \(x\) in A. We have \(\chi(jx) = \chi(x)\) for all \(\chi \in \mathsf{X}(\mathsf{A})\) , hence \(jx = x\) since A is without radical. Thus, \(j\) is an identity element of A.

PROPOSITION 15. --- Let A be a commutative Banach algebra, let \(I_{1}\) be an ideal of A and \(F_{1}\) be the set of \(\chi \in \mathsf{X}(A)\) which vanish on \(I_{1}\) . Let \(F_{2}\) be a subset of X(A) disjoint from \(F_{1}\) , closed for the Jacobson topology, and compact for the weak topology. Then there exists \(u \in \text{I}_{1}\) such that \(\mathcal{G}(u) = 1\) on \(F_{2}\) .

Let \(I_{2}\) be the intersection of the kernels of the characters belonging to \(F_{2}\) . The Banach algebra A/ \(I_{2}\) is without radical (prop. 8 of I, p.~38). Since \(F_{2}\) is closed for the Jacobson topology, the only elements of X(A) vanishing on \(I_{2}\) are those of \(F_{2}\) (cf.~I, p.~13). Hence \(F_{2}\) , equipped with the topology induced by the weak topology of X(A), is identified with X(A/ \(I_{2}\) ) equipped with the weak topology (I, p.~9, no. 7). Since \(F_{2}\) is weakly compact, the algebra A/ \(I_{2}\) possesses an identity element (prop. 14).

We then have \(I_{1} + I_{2} = A\) . Indeed, in the contrary case, \((I_{1} + I_{2})/I_{2}\) would be a proper ideal, hence contained in the kernel of a nonzero character of \(A/I_{2}\) (I, p.~30, th. 2). This character would define, by composition with the canonical projection \(A \rightarrow A/I_{2}\) , a nonzero character \(\chi\) of A vanishing on \(I_{1}\) and \(I_{2}\) , and would therefore belong to \(F_{1} \cap F_{2}\) , contrary to the hypothesis.

Since \(I_{1} + I_{2} = A\) , there exists \(u \in I_{1}\) whose class in \(A/I_{2}\) is an identity element of \(A/I_{2}\) . Then \(\chi(u) = 1\) for all \(\chi \in F_{2}\) , which concludes the proof.

COROLLARY. --- Let A be a commutative Banach algebra. Let \(F_{1}\) and \(F_{2}\) two disjoint subsets of \(\mathsf{X}(\mathsf{A})\) , closed for the Jacobson topology. We assume \(F_{2}\) weakly compact. Then there exists \(u \in A\) such that \(\mathcal{G}(u) = 1\) on \(F_{2}\) and \(\mathcal{G}(u) = 0\) on \(F_{1}\) .

